\documentclass[12pt]{article}
\usepackage[utf8]{inputenc}
\usepackage[T1]{fontenc}
\usepackage{hyperref}
\title{Puare Language Resources for K-12 Curriculum Development}
\author{Compiled by Claude Sonnet 5.5 (claude-sonnet-5-5) for Andr\'as Kornai}
\date{}
\begin{document}
\maketitle

\section{Introduction}
Puare (also called Puari; ISO 639-3 \texttt{pux}) is a Papuan language of Sandaun Province, Papua New Guinea, classified in the Sko (Skou) family. It is critically endangered. Search results report about 35 speakers in 2003 (Joshua Project lists 60), and a 1998 tsunami reportedly killed most speakers. No curriculum materials, primers or school programmes were found.

\section{Primary Educational Resources}
No primers, textbooks or school materials in or about Puare were found in search results.

\section{Linguistic and Reference Materials}
\begin{itemize}
\item Wikipedia, Puari language: overview of status, classification and sources. \url{https://en.wikipedia.org/wiki/Puari_language}
\item Wikipedia, Skou languages: family overview and the classification debate. \url{https://en.wikipedia.org/wiki/Skou_languages}
\item ASJP database entry for Puare (a word list). \url{https://asjp.clld.org/languages/PUARE}
\end{itemize}

\section{Religious and Community Resources}
\begin{itemize}
\item Joshua Project, Puari in Papua New Guinea: people-group profile. \url{https://joshuaproject.net/people_groups/14464}
\end{itemize}
No scripture translations or community organisations were found.

\section{Digital and Audio Resources}
\begin{itemize}
\item Endangered Languages Project entry for Puare. \url{https://www.endangeredlanguages.com/lang/2647?hl=en}
\end{itemize}
Search results state that PARADISEC holds Puare tape recordings made by Don Laycock and Melissa Crowther. I did not see a direct collection URL, so none is given.

\section{Curriculum Development Context}
Documentation is very thin and the speaker community is tiny. Any curriculum would have to start with documentation work with the remaining speakers and the archived recordings, not with existing teaching materials.

\end{document}